\textbf{Phần 1. Trắc nghiệm lựa chọn}
\setcounterpageref{ex}{0}

\begin{ex} %[4D8N48]
Kết quả của phép nhân $450 \times 20$ là
\choice
{$900$}
{$90\,000$}
{$900\,000$}
{\True $9000$}
\loigiai{
$450 \times 20 = 9000$.\\
\textbf{Đáp án của câu hỏi là D}
}
\end{ex}

\begin{ex} %[4D8N48]
Kết quả của phép chia $3600 : 60$ là
\choice
{\True $60$}
{$6$}
{$600$}
{$36$}
\loigiai{
$3600 : 60 = 60$.\\
\textbf{Đáp án của câu hỏi là A}
}
\end{ex}

\begin{ex} %[4D8H48]
Tính thuận tiện: $25 \times 78 \times 4$.
\choice
{$7600$}
{\True $7800$}
{$8000$}
{$7500$}
\loigiai{
$(25 \times 4) \times 78 = 100 \times 78 = 7800$.\\
\textbf{Đáp án của câu hỏi là B}
}
\end{ex}

\begin{ex} %[4D8H48]
Tính thuận tiện: $35 \times 46 + 35 \times 54$.
\choice
{\True $3500$}
{$3400$}
{$3600$}
{$3800$}
\loigiai{
$35 \times (46 + 54) = 35 \times 100 = 3500$.\\
\textbf{Đáp án của câu hỏi là A}
}
\end{ex}

\begin{ex} %[4D8H48]
Trung bình cộng của ba số: $40, 60, 80$ là
\choice
{$50$}
{$70$}
{\True $60$}
{$180$}
\loigiai{
$(40 + 60 + 80) : 3 = 60$.\\
\textbf{Đáp án của câu hỏi là C}
}
\end{ex}

\begin{ex} %[4D8H48]
Mua 4 chiếc bút hết $32\,000$ đồng. Mua 9 chiếc bút như thế hết số tiền là
\choice
{$64\,000$ đồng}
{\True $72\,000$ đồng}
{$80\,000$ đồng}
{$76\,000$ đồng}
\loigiai{
Giá 1 chiếc: $32\,000 : 4 = 8000$ đồng.\\
9 chiếc: $8000 \times 9 = 72\,000$ đồng.\\
\textbf{Đáp án của câu hỏi là B}
}
\end{ex}

\begin{ex} %[4D8V48]
Một mảnh đất hình chữ nhật có chiều dài $35\text{ m}$, chiều rộng $20\text{ m}$. Diện tích mảnh đất là
\choice
{$600\text{ m}^2$}
{$800\text{ m}^2$}
{$110\text{ m}^2$}
{\True $700\text{ m}^2$}
\loigiai{
$35 \times 20 = 700\text{ m}^2$.\\
\textbf{Đáp án của câu hỏi là D}
}
\end{ex}

\begin{ex} %[4D8V48]
Một ô tô trong 3 giờ đầu mỗi giờ đi được $45\text{ km}$, trong 2 giờ sau mỗi giờ đi được $55\text{ km}$. Trung bình mỗi giờ ô tô đi được bao nhiêu ki-lô-mét?
\choice
{$48\text{ km}$}
{\True $49\text{ km}$}
{$50\text{ km}$}
{$52\text{ km}$}
\loigiai{
Tổng quãng đường: $3 \times 45 + 2 \times 55 = 135 + 110 = 245\text{ km}$.\\
Trung bình: $245 : 5 = 49\text{ km}$.\\
\textbf{Đáp án của câu hỏi là B}
}
\end{ex}

\begin{ex} %[4D8V48]
Tìm $x$ biết: $x \times 45 = 1800$.
\choice
{$30$}
{$50$}
{\True $40$}
{$60$}
\loigiai{
$x = 1800 : 45 = 40$.\\
\textbf{Đáp án của câu hỏi là C}
}
\end{ex}

\begin{ex} %[4D8V48]
Tính nhanh giá trị biểu thức: $125 \times 16$.
\choice
{\True $2000$}
{$1800$}
{$2200$}
{$2400$}
\loigiai{
$125 \times 8 \times 2 = 1000 \times 2 = 2000$.\\
\textbf{Đáp án của câu hỏi là A}
}
\end{ex}

\begin{ex} %[4D8V48]
Một đội công nhân nhận sửa $1200\text{ m}$ đường. Trong 5 ngày đầu sửa được $600\text{ m}$. Hỏi với năng suất đó, đội sửa xong toàn bộ đoạn đường trong bao nhiêu ngày?
\choice
{$8$ ngày}
{$12$ ngày}
{$15$ ngày}
{\True $10$ ngày}
\loigiai{
Mỗi ngày sửa: $600 : 5 = 120\text{ m}$.\\
Toàn bộ đoạn đường: $1200 : 120 = 10$ ngày.\\
\textbf{Đáp án của câu hỏi là D}
}
\end{ex}

\begin{ex} %[4D8C48]
Tính giá trị biểu thức: $A = (1 + 3 + 5 + \dots + 99) \times (45 \times 2 - 90)$.
\choice
{$2500$}
{$9900$}
{\True $0$}
{$10\,000$}
\loigiai{
Vì $45 \times 2 - 90 = 90 - 90 = 0$ nên tích $A = 0$.\\
\textbf{Đáp án của câu hỏi là C}
}
\end{ex}

\textbf{Phần 2. Trắc nghiệm đúng sai}
\setcounterpageref{ex}{0}

\begin{ex} %[4D8H48]
Xét tính đúng sai của các biểu thức phép nhân chia:
\choiceTF
{\True $a \times (b + c) = a \times b + a \times c$.}
{\True $125 \times 8 = 1000$.}
{$250 \times 40 = 1000$.}
{\True $250 \times 40 = 10\,000$.}
\loigiai{
\begin{itemize}
    \item Tính chất phân phối: Đúng.
    \item $125 \times 8 = 1000$. Đúng.
    \item $250 \times 40 = 10\,000 \neq 1000$. Ý 3 Sai, Ý 4 Đúng.
\end{itemize}
\begin{itemize}
    \item a) \textbf{Đúng}.
    \item b) \textbf{Đúng}.
    \item c) \textbf{Sai}.
    \item d) \textbf{Đúng}.
\end{itemize}
}
\end{ex}

\begin{ex} %[4D8V48]
Hai tổ sản xuất khẩu trang: Tổ 1 có 8 người may được 640 hộp, tổ 2 có 10 người may được 800 hộp. Xét tính đúng sai:
\choiceTF
{\True Trung bình mỗi người tổ 1 may được 80 hộp.}
{\True Trung bình mỗi người tổ 2 may được 80 hộp.}
{Năng suất lao động trung bình của mỗi người ở tổ 2 cao hơn tổ 1.}
{\True Cả hai tổ có mức năng suất trung bình của mỗi người như nhau.}
\loigiai{
\begin{itemize}
    \item Tổ 1: $640 : 8 = 80$ hộp. Đúng.
    \item Tổ 2: $800 : 10 = 80$ hộp. Đúng.
    \item Năng suất bằng nhau (80 hộp/người). Ý 3 Sai, Ý 4 Đúng.
\end{itemize}
\begin{itemize}
    \item a) \textbf{Đúng}.
    \item b) \textbf{Đúng}.
    \item c) \textbf{Sai}.
    \item d) \textbf{Đúng}.
\end{itemize}
}
\end{ex}

\textbf{Phần 3. Trả lời ngắn}
\setcounterpageref{ex}{0}

\begin{ex} %[4D8H48]
Tính thuận tiện: $25 \times 64 \times 4$.
\shortans{6400}
\loigiai{
\textbf{Đáp số:} $6400$\\
$(25 \times 4) \times 64 = 6400$.
}
\end{ex}

\begin{ex} %[4D8H48]
Tìm số trung bình cộng của 3 số: $50, 70, 90$.
\shortans{70}
\loigiai{
\textbf{Đáp số:} $70$\\
$(50 + 70 + 90) : 3 = 70$.
}
\end{ex}

\begin{ex} %[4D8V48]
Mua 6 quyển vở hết $48\,000$ đồng. Hỏi mua 10 quyển vở như thế hết bao nhiêu nghìn đồng (chỉ ghi số nghìn đồng)?
\shortans{80}
\loigiai{
\textbf{Đáp số:} $80$\\
$48 : 6 \times 10 = 80$ nghìn đồng.
}
\end{ex}

\begin{ex} %[4D8V48]
Tính giá trị biểu thức: $45 \times 101 - 45$.
\shortans{4500}
\loigiai{
\textbf{Đáp số:} $4500$\\
$45 \times 100 = 4500$.
}
\end{ex}

\begin{ex} %[4D8V48]
Tìm $x$ biết: $x : 35 = 200$.
\shortans{7000}
\loigiai{
\textbf{Đáp số:} $7000$\\
$x = 200 \times 35 = 7000$.
}
\end{ex}

\begin{ex} %[4D8C48]
Tính tích: $A = 25 \times 4 \times 125 \times 8 \times 50 \times 2$.
\shortans{10000000}
\loigiai{
\textbf{Đáp số:} $10000000$\\
$100 \times 1000 \times 100 = 10\,000\,000$.
}
\end{ex}
